% !TeX root = ../navier-stokes-manuscript.tex
% ============================================================
% 1.7 Singularity and the Derivative Tower
% ============================================================
\subsection{Singularity and the Derivative Tower}\label{sub:SingularityAndTheDerivativeTower}

A finite limiting velocity along one material parcel does not settle whether the velocity field extends smoothly to that time.
Let \(X(a,t)\) be the trajectory of the parcel labelled by \(a\), and write its carried velocity as
\[
  \dot X(a,t)=u(t,X(a,t)),
  \qquad
  U(a,t):=u(t,X(a,t)).
\]
The curve \(U(a,\cdot)\) is the velocity that the parcel actually experiences.

A simple time profile shows how much its terminal value can hide.
Fix a unit vector \(e\) and suppose, kinematically, that near a finite time \(T\),
\[
  U(t)=U_*+(T-t)^{3/2}e,
  \qquad t<T.
\]
Then
\[
  U(t)\longrightarrow U_*,
  \qquad
  \dot U(t)=-\frac32(T-t)^{1/2}e\longrightarrow0,
  \qquad
  |\ddot U(t)|=\frac34(T-t)^{-1/2}|e|\longrightarrow\infty.
\]
The velocity arrives, and even the acceleration tends to zero, while the jerk diverges.
This is a kinematic profile, with no claim that a \NavierStokes{} solution realizes it.
It isolates one kind of singularity: lower responses can arrive at \(T\) while a higher response does not.

The \NavierStokes{} equation already contains the parcel's first response.
Set
\[
  D_t:=\partial_t+u\cdot\nabla.
\]
The chain rule along \(X(a,t)\) and the momentum equation give
\[
  \dot U(a,t)
  =
  D_tu(t,X(a,t))
  =
  \bigl(-\nabla p+\nu\Delta u\bigr)(t,X(a,t)).
\]
The pressure force and the viscous response to spatial curvature determine how that same parcel accelerates.
Differentiate the same equality along the same path, and its jerk is
\[
  \ddot U(a,t)
  =
  D_t^2u(t,X(a,t))
  =
  D_t\bigl(-\nabla p+\nu\Delta u\bigr)(t,X(a,t)).
\]
Another material derivative gives the rate at which the jerk changes, and the process continues at every finite order while the field remains smooth.
Higher names the derivative order.
The response may shrink, pass through zero, or vanish.
Each material rung is the derivative of the one below, evaluated from the same \NavierStokes{} field.

Material differentiation follows the parcel, while Eulerian differentiation watches the field at one fixed spatial coordinate.
The identity \(D_t=\partial_t+u\cdot\nabla\) joins the two views: the parcel changes in time while it also moves through spatial variation.
To see that joining inside the equation, write the fixed-coordinate derivatives as
\[
  V_n:=\partial_t^nu,
  \qquad
  Q_n:=\partial_t^np,
  \qquad n\geq0.
\]
The first Eulerian row is
\[
  V_1
  +(V_0\cdot\nabla)V_0
  +\nabla Q_0
  =
  \nu\Delta V_0.
\]
Differentiate once in time and the next row becomes
\[
  V_2
  +(V_1\cdot\nabla)V_0
  +(V_0\cdot\nabla)V_1
  +\nabla Q_1
  =
  \nu\Delta V_1.
\]
The time derivative can land on either velocity factor.
On the classical interval, fixed-coordinate time derivatives commute with \(\nabla\) and \(\Delta\), so the pressure gradient and viscous Laplacian reappear at the new time order.
Repeating this Eulerian differentiation produces the higher temporal rows.
Differentiating those identities in space produces the spatial and mixed rows.
Every resulting row still couples temporal response, spatial variation, viscous curvature, and pressure.

Using a multi-index \(\alpha\) for the spatial derivatives, collect the velocity coordinates at one spacetime point in
\[
  \mathcal J^\infty u(t,x)
  :=
  \bigl(
    \partial_t^n\partial_x^\alpha u(t,x)
  \bigr)_{n\geq0,\,\alpha\in\mathbb N_0^3}.
\]
This is the derivative tower of the velocity.
At each classical instant, its rungs are simultaneous readings of one velocity field.
Material acceleration, jerk, and every higher carried response are combinations of these fixed-coordinate derivatives along the trajectory.
The pressure rows \(Q_n\) stay attached through the differentiated momentum equations and the divergence-free constraint.

A finite terminal value of \(U\) leaves the higher tower unsettled.
Every required temporal, spatial, and mixed derivative still has to reach \(T\) as part of the same smooth field.
Singularity can sit at any rung where that extension fails.
